\bookchapter{Kapitel 10: Mobiles Bezahlen und Alltagsdienste}

\chapterintro{Ganz ehrlich: Meine Eltern und deren Generation waren gegenüber dem „Bezahlen mit dem Smartphone“ anfangs skeptisch. „Das Geld ist im Smartphone, was, wenn es gestohlen wird?” Diese Sorge ist verständlich, aber die Realität ist: Die Wahrscheinlichkeit, dass Bargeld in Ihrer Hosentasche verloren geht, ist viel höher als die, dass das mobile Bezahlen missbraucht wird. Mobiles Bezahlen bedeutet nicht, dass das Geld im Smartphone aufbewahrt wird – es ist nur ein Schlüssel, ein Schlüssel, der Ihr Bankkonto öffnet.}

\booksection{Wie bindet man eigentlich eine Bankkarte ein?}

Egal ob Sie Alipay oder WeChat verwenden, beim ersten Mal die Karte zu hinterlegen, sind diese Schritte nötig:

\begin{enumerate}
  \item Öffnen Sie die App, finden Sie „Mein" „Geldbörse". Tippen Sie auf „Bankkarte hinzufügen".
  \item Kartennummer eingeben. Die meisten Apps unterstützen jetzt das direkte Fotografieren zum Erkennen der Kartennummer, ohne sie Ziffer für Ziffer einzutippen.
  \item Die Bank sendet einen Bestätigungscode per SMS an die hinterlegte Handynummer. Tragen Sie ihn ein und legen Sie ein 6-stelliges Zahlungs-Passwort fest.
  \item Dieses Zahlungs-Passwort, das Passwort Ihrer Bankkarte und das Sperrbildschirm-Passwort sollten am besten drei verschiedene Passwörter sein – sparen Sie nicht an Bequemlichkeit, indem Sie alle gleich setzen.
\end{enumerate}

Ein kleiner Tipp: Wenn Sie sich Sorgen um die Sicherheit machen, können Sie in der Bank-App ein elektronisches Konto der Klasse II oder III eröffnen, das ausschließlich für mobiles Bezahlen hinterlegt wird. Auf dieser Karte liegt nur das Geld für den täglichen Bedarf – selbst wenn etwas passiert, ist der Schaden begrenzt.

\booksection{Zahlungscode: Wer sieht, wer bekommt das Geld?}

Viele Menschen verwechseln Zahlungscode und Empfangscode – dieser Unterschied ist wichtig.

\begin{itemize}
  \item \textbf{Zahlungscode}: Das ist der Strichcode/QR-Code, den Sie dem Kassierer zum Scannen zeigen. Dieser Code entspricht Ihrer „elektronischen Geldbörse" – wer ihn scannt, kann Geld von Ihrem Konto abziehen.\textbf{Also niemals einen Screenshot an jemanden schicken!}Dieser Code wird jede Minute automatisch aktualisiert, um Missbrauch zu verhindern.
  \item \textbf{Empfangscode}: Jemand scannt diesen Code, um Ihnen Geld zu überweisen. Dieser Code kann als Screenshot gespeichert und an andere geschickt werden – sicher.
\end{itemize}

Wenn Sie diesen Unterschied verinnerlichen, können Sie viele Betrugsmaschen umgehen. Wer immer Sie auffordert, „den Zahlungscode als Screenshot zu schicken“ – zu 100 \% Betrüger.

\booksection{Diese Sicherheitseinstellungen – 5 Minuten investieren, viel Ärger sparen}

\begin{itemize}
  \item \textbf{Fingerabdruck/Gesichtszahlung}: Aktivieren Sie in den Zahlungseinstellungen von Alipay und WeChat die biometrische Zahlung. So ist bei jeder Zahlung ein Fingerabdruck oder Gesichtsscan nötig – selbst wenn das Smartphone verloren geht, kann der Finder nicht Ihr Geld ausgeben.
  \item \textbf{Sicherheitssperre}: Sowohl Alipay als auch WeChat haben eine „Sicherheitssperre"-Funktion. Wenn aktiviert, ist beim Öffnen der „Geldbörse"-Seite der App eine zusätzliche Identitätsprüfung nötig. Pfad: Alipay Mein Einstellungen Konto und Sicherheit Sicherheitscenter Sicherheitssperre.
  \item \textbf{Prüfung der passwortfreien Zahlung}: Viele haben verschiedene passwortfreie Zahlungen eingerichtet und vergessen. Gehen Sie in die „Zahlungseinstellungen" von Alipay und WeChat, suchen Sie „Passwortfreie Zahlung / Automatischer Einzug" und prüfen Sie, welche Dienste aktiviert sind – schalten Sie nicht mehr benötigte ab.
  \item \textbf{Verwaltung der Anmeldegeräte}: Prüfen Sie regelmäßig, welche Geräte sich in Ihrem Konto angemeldet haben. Wenn Sie ein unbekanntes Gerät entdecken, melden Sie es sofort ab. Pfad: Alipay Einstellungen Konto und Sicherheit Verwaltung der Anmeldegeräte.
\end{itemize}

\booksection{Bus und U-Bahn: Das Smartphone ist die Fahrkarte}

Heute nutzen beim Bus- und U-Bahn-Fahren kaum noch jemand physische Fahrkarten. Für das Smartphone an der Sperre gibt es zwei Wege:

\begin{itemize}
  \item \textbf{Virtuelle Fahrkarte in der Smartphone-Geldbörse}: In Apple Wallet, Huawei Wallet oder Xiaomi Wallet lässt sich direkt eine virtuelle Fahrkarte aktivieren. Der Vorteil: Beim Entwerten muss das Smartphone nicht entsperrt werden, einfach anhalten genügt, die Reaktionsgeschwindigkeit ist am höchsten, und es funktioniert auch offline. In manchen Städten bietet die virtuelle Fahrkarte die gleichen Umsteigevorteile wie die physische.
  \item \textbf{Fahrcode von Alipay / WeChat}: Suchen Sie in der App nach „Fahrcode" und aktivieren Sie den Bus- oder U-Bahn-Code Ihrer Stadt. Der Vorteil: Keine Aufladung nötig, es wird direkt vom Guthaben oder von der hinterlegten Bankkarte abgebucht. Aber es ist eine Netzverbindung nötig, und vor dem Scannen muss die App geöffnet werden, um den Code zu finden.
\end{itemize}

Ich persönlich empfehle die virtuelle Fahrkarte in der Smartphone-Geldbörse. Auch wenn man beim ersten Mal 10–20 Yuan aufladen muss, ist die Nutzung wesentlich bequemer, besonders wenn man zur U-Bahn hetzt – man muss nicht anhalten, um das Smartphone zu entsperren und die App zu suchen.

\booksection{Alltagsgebühren, Essenslieferung, Fahrdienste}

\begin{itemize}
  \item \textbf{Wasser, Strom, Gas}: In „Alltagsgebühren" von Alipay oder WeChat hinterlegen Sie Ihre Zählernummern für Wasser, Strom und Gas. Wenn jeden Monat die Rechnung kommt, tippen Sie einfach hinein und zahlen. Die Zählernummer finden Sie meist auf der Papierrechnung zu Hause oder erfragen Sie beim Kundenservice.
  \item \textbf{Essenslieferung vergleichen}: Dasselbe Restaurant hat auf Meituan und Ele.me oft unterschiedliche Preise und Angebote. Schauen Sie vor der Bestellung auf beiden Plattformen, manchmal sparen Sie ein paar Yuan. Außerdem gibt es bei Meituan die Funktion „Gemeinsam bestellen" – wenn Sie mit Freunden zusammen bestellen, wird es deutlich günstiger.
  \item \textbf{Sicherheit bei Fahrdiensten}: In DiDi oder Gaode Ride sollten Sie „Notfallkontakt" und „Fahrt teilen" einrichten. Während der Fahrt teilt die App Ihre Fahrt automatisch mit den hinterlegten Kontakten. Außerdem: Bei nächtlichen Fahrten bevorzugen Sie „Taxi" statt „Schnellwagen" – Taxis unterliegen einer Branchenaufsicht und haben eine höhere Sicherheit.
\end{itemize}

\booksection{Krankenversicherung, Sozialversicherung und Behördengänge mit dem Smartphone}

In den letzten Jahren hat sich viel verändert – viele Dinge, für die man früher in Schlangen vor Schaltern stand, sind heute in wenigen Minuten auf dem Smartphone erledigt.

\begin{itemize}
  \item \textbf{Elektronische Krankenversicherung}: Suchen Sie in Alipay oder WeChat nach „Elektronische Krankenversicherung". Nach der Aktivierung reicht beim Arztbesuch oder in der Apotheke das Smartphone – keine physische Karte mehr nötig. Beim ersten Mal ist eine Gesichtserkennung nötig, danach wird bei jeder Nutzung ein dynamischer Code generiert.
  \item \textbf{Sozialversicherung abfragen}: Laden Sie die App „Zhangshang 12333" herunter oder nutzen Sie in Alipay das „Bürgercenter", um Beitragsaufzeichnungen, Rentenkontostand und Arbeitslosengeld abzufragen.
  \item \textbf{Einkommensteuer}: Die jährliche Steuererklärung von März bis Juni lässt sich in der App „Persönliche Einkommensteuer" erledigen. Prüfen Sie, ob Sie Sonderabzüge geltend machen können (Kinderbildung, Kreditzinsen, Unterhalt für Ältere usw.) – diese können Ihre Steuerlast mindern oder zu einer Rückerstattung führen.
\end{itemize}

\begin{tcbwarn}{Merken Sie sich diese Punkte, um Betrug zu vermeiden}
\begin{itemize}
  \item Polizei, Staatsanwaltschaft, Banken oder Kundenservice werden Sie niemals auffordern, auf ein „Sicherheitskonto" zu überweisen.
  \item Wer auch immer Sie auffordert, den Zahlungscode als Screenshot zu schicken – Betrüger.
  \item Klicken Sie nicht auf unbekannte Links, besonders bei SMS wie „ETC abgelaufen" oder „Sozialversicherungskarte gesperrt".
  \item Rückerstattungen bei Online-Einkäufen werden immer auf dem ursprünglichen Zahlungsweg zurückgebucht – Sie müssen keine Bankkartennummer angeben.
\end{itemize}

\end{tcbwarn}

\begin{tcbpractice}{Jetzt tun: Führen Sie eine Sicherheitsüberprüfung Ihrer Zahlungs-Apps durch}
\begin{enumerate}
  \item Öffnen Sie Alipay und WeChat und bestätigen Sie, dass Fingerabdruck-/Gesichtszahlung aktiviert ist.
  \item Finden Sie die Funktion „Sicherheitssperre" und aktivieren Sie sie.
  \item Prüfen Sie die Liste „Passwortfreie Zahlung / Automatischer Einzug" und schalten Sie nicht mehr benötigte Dienste ab.
  \item Sehen Sie die Liste „Anmeldegeräte" an und entfernen Sie unbekannte Geräte.
  \item Aktivieren Sie Ihre elektronische Krankenversicherung – beim nächsten Arztbesuch brauchen Sie keine physische Karte mehr.
\end{enumerate}

\end{tcbpractice}
